\textbf{Phần 1. Trắc nghiệm lựa chọn}
\setcounterpageref{ex}{0}

\begin{ex} %[4D1V01]
Số $X$ là số có 5 chữ số. Biết chữ số hàng chục nghìn là số lớn nhất có một chữ số, chữ số hàng trăm kém chữ số hàng chục nghìn $4$ đơn vị, chữ số hàng đơn vị là $3$, các chữ số còn lại đều bằng $0$. Cách đọc số $X$ là
\choice
{Chín mươi nghìn năm trăm ba mươi}
{Chín mươi lăm nghìn không trăm linh ba}
{\True Chín mươi nghìn năm trăm linh ba}
{Chín mươi nghìn năm trăm linh năm}
\loigiai{
Chữ số hàng chục nghìn là 9.\\
Chữ số hàng trăm là: $9 - 4 = 5$.\\
Chữ số hàng đơn vị là 3.\\
Các chữ số hàng nghìn và hàng chục đều bằng 0.\\
Do đó, số $X$ được viết là $90\,503$ và đọc là: \textit{Chín mươi nghìn năm trăm linh ba}.\\
\textbf{Đáp án của câu hỏi là C}
}
\end{ex}

\begin{ex} %[4D1V01]
Trong số $84\,542$, giá trị của chữ số $4$ ở hàng nghìn gấp giá trị của chữ số $4$ ở hàng chục số lần là
\choice
{\True $100$ lần}
{$10$ lần}
{$1000$ lần}
{$400$ lần}
\loigiai{
Chữ số 4 ở hàng nghìn có giá trị là $4000$.\\
Chữ số 4 ở hàng chục có giá trị là $40$.\\
Giá trị của chữ số 4 ở hàng nghìn gấp giá trị của chữ số 4 ở hàng chục số lần là:
$$4000 : 40 = 100\text{ (lần)}$$
\textbf{Đáp án của câu hỏi là A}
}
\end{ex}

\begin{ex} %[4D1V01]
Số liền sau của số chẵn lớn nhất có 5 chữ số khác nhau là
\choice
{$98\,764$}
{$98\,767$}
{$99\,999$}
{\True $98\,765$}
\loigiai{
Số chẵn lớn nhất có 5 chữ số khác nhau là $98\,764$ (các chữ số được chọn giảm dần là 9, 8, 7, 6 và chữ số hàng đơn vị chẵn là 4).\\
Số liền sau của $98\,764$ là: $98\,764 + 1 = 98\,765$.\\
\textbf{Đáp án của câu hỏi là D}
}
\end{ex}

\begin{ex} %[4D1V01]
Số lẻ bé nhất có 5 chữ số khác nhau mà chữ số hàng chục nghìn bằng $3$ là
\choice
{$30\,215$}
{\True $30\,125$}
{$31\,025$}
{$30\,145$}
\loigiai{
Để số là bé nhất, ta chọn các chữ số từ trái sang phải bé nhất có thể:\\
- Chữ số hàng chục nghìn là 3 (theo đề bài).\\
- Chữ số hàng nghìn bé nhất khác 3 là 0.\\
- Chữ số hàng trăm bé nhất khác 3 và 0 là 1.\\
- Chữ số hàng chục bé nhất khác 3, 0 và 1 là 2.\\
- Chữ số hàng đơn vị phải là số lẻ khác 3 và 1 nên chọn chữ số 5.\\
Vậy số cần tìm là $30\,125$.\\
\textbf{Đáp án của câu hỏi là B}
}
\end{ex}

\begin{ex} %[4D1V01]
Một số gồm $5$ chục nghìn, $24$ trăm và $7$ đơn vị. Số đó được viết là
\choice
{$52\,470$}
{$50\,247$}
{\True $52\,407$}
{$52\,047$}
\loigiai{
Ta có:\\
- $5$ chục nghìn $= 50\,000$\\
- $24$ trăm $= 2400$\\
- $7$ đơn vị $= 7$\\
Số đó là: $50\,000 + 2400 + 7 = 52\,407$.\\
\textbf{Đáp án của câu hỏi là C}
}
\end{ex}

\begin{ex} %[4D1V01]
Cho các biểu thức và số sau:
$M = 40\,000 + 6000 + 300 + 20$;
$N = 46\,350$;
$P = 46\,000 + 35$;
$Q$ là số liền sau của số $46\,309$.
Sắp xếp các số $M, N, P, Q$ theo thứ tự từ lớn đến bé là
\choice
{\True $N > M > Q > P$}
{$N > Q > M > P$}
{$M > N > Q > P$}
{$N > M > P > Q$}
\loigiai{
Ta tính giá trị của từng số:\\
- $M = 46\,320$\\
- $N = 46\,350$\\
- $P = 46\,035$\\
- $Q = 46\,309 + 1 = 46\,310$\\
Vì $46\,350 > 46\,320 > 46\,310 > 46\,035$ nên thứ tự từ lớn đến bé là: $N > M > Q > P$.\\
\textbf{Đáp án của câu hỏi là A}
}
\end{ex}

\begin{ex} %[4D1V01]
Cho số $\overline{6a,725}$. Biết rằng khi làm tròn số này đến hàng chục nghìn ta được số $70\,000$. Chữ số $a$ bé nhất có thể nhận là
\choice
{$4$}
{$6$}
{$0$}
{\True $5$}
\loigiai{
Để làm tròn số $\overline{6a,725}$ đến hàng chục nghìn được $70\,000$ (tức là làm tròn lên), chữ số hàng nghìn $a$ phải thỏa mãn điều kiện $a \ge 5$.\\
Do đó, chữ số $a$ bé nhất có thể nhận là $5$.\\
\textbf{Đáp án của câu hỏi là D}
}
\end{ex}

\begin{ex} %[4D1V01]
Tìm số tự nhiên có 5 chữ số khác nhau, biết chữ số hàng chục nghìn là $6$, các chữ số tiếp theo giảm dần từ trái sang phải và chữ số hàng đơn vị là một số chẵn lớn hơn 0. Số lớn nhất thỏa mãn điều kiện đó là
\choice
{$65\,420$}
{\True $65\,432$}
{$65\,430$}
{$65\,320$}
\loigiai{
Gọi số cần tìm có dạng $\overline{6abcd}$ với $6 > a > b > c > d$.\\
Vì chữ số hàng đơn vị $d$ là số chẵn lớn hơn 0 nên $d \ge 2$.\\
Giữa 6 và 2 có đúng ba chữ số là 5; 4; 3. Vì các chữ số giảm dần nên ta có: $a = 5, b = 4, c = 3$ và $d = 2$.\\
Vậy số cần tìm là $65\,432$.\\
\textbf{Đáp án của câu hỏi là B}
}
\end{ex}

\begin{ex} %[4D1V01]
Cho 4 chữ số: $0; 2; 5; 7$. Có thể viết được bao nhiêu số có 4 chữ số khác nhau và chia hết cho $5$ từ các chữ số đã cho?
\choice
{$18$ số}
{$12$ số}
{\True $10$ số}
{$6$ số}
\loigiai{
Một số chia hết cho 5 khi chữ số hàng đơn vị là 0 hoặc 5.\\
\textbf{Trường hợp 1:} Chữ số hàng đơn vị là 0.\\
- Hàng nghìn có 3 cách chọn (từ các chữ số 2; 5; 7).\\
- Hàng trăm có 2 cách chọn.\\
- Hàng chục có 1 cách chọn.\\
Số các số tạo thành là: $3 \times 2 \times 1 = 6$ (số).\\
\textbf{Trường hợp 2:} Chữ số hàng đơn vị là 5.\\
- Hàng nghìn có 2 cách chọn (chọn 2 hoặc 7, vì phải khác 0).\\
- Hàng trăm có 2 cách chọn (chọn trong các chữ số còn lại bao gồm cả 0).\\
- Hàng chục có 1 cách chọn.\\
Số các số tạo thành là: $2 \times 2 \times 1 = 4$ (số).\\
Tổng số các số viết được là: $6 + 4 = 10$ (số).\\
\textbf{Đáp án của câu hỏi là C}
}
\end{ex}

\begin{ex} %[4D1V01]
Có bao nhiêu số tự nhiên có 5 chữ số mà khi làm tròn đến hàng chục nghìn ta được số $80\,000$?
\choice
{\True $10\,000$ số}
{$9999$ số}
{$5000$ số}
{$9000$ số}
\loigiai{
Các số tự nhiên có 5 chữ số khi làm tròn đến hàng chục nghìn được $80\,000$ là các số từ $75\,000$ đến $84\,999$.\\
Số lượng các số thỏa mãn là:
$$84\,999 - 75\,000 + 1 = 10\,000\text{ (số)}$$
\textbf{Đáp án của câu hỏi là A}
}
\end{ex}

\begin{ex} %[4D1V01]
Một kho lương thực có $28\,400$ kg gạo. Tuần thứ nhất kho xuất đi $6200$ kg gạo, tuần thứ hai kho xuất đi số gạo gấp đôi tuần thứ nhất. Số ki-lô-gam gạo còn lại trong kho sau hai tuần khi làm tròn đến hàng nghìn là
\choice
{$9000$ kg}
{$16\,000$ kg}
{$12\,000$ kg}
{\True $10\,000$ kg}
\loigiai{
Tuần thứ hai kho xuất đi số gạo là:
$$6200 \times 2 = 12\,400\text{ (kg)}$$
Cả hai tuần kho đã xuất đi tổng số gạo là:
$$6200 + 12\,400 = 18\,600\text{ (kg)}$$
Số gạo còn lại trong kho sau hai tuần là:
$$28\,400 - 18\,600 = 9800\text{ (kg)}$$
Làm tròn số $9800$ đến hàng nghìn (vì chữ số hàng trăm là 8 nên làm tròn lên), ta được $10\,000$ kg.\\
\textbf{Đáp án của câu hỏi là D}
}
\end{ex}

\begin{ex} %[4D1V01]
Viết số lớn nhất có 5 chữ số khác nhau mà tổng các chữ số của số đó bằng $18$ và số đó chia hết cho $2$.
\choice
{$95\,310$}
{\True $96\,210$}
{$96\,102$}
{$85\,320$}
\loigiai{
Để số là lớn nhất, ta chọn chữ số hàng chục nghìn lớn nhất là 9.\\
Tổng của 4 chữ số còn lại là: $18 - 9 = 9$.\\
Số chia hết cho 2 nên chữ số hàng đơn vị phải là số chẵn: chọn chữ số hàng đơn vị là 0 để số này nhỏ nhất có thể ở hàng cuối và là số chẵn.\\
Khi đó tổng 3 chữ số hàng nghìn, hàng trăm và hàng chục bằng 9.\\
Để số lớn nhất, chữ số hàng nghìn phải lớn nhất có thể:\\
- Nếu hàng nghìn là 7 thì hai chữ số còn lại có tổng bằng $9 - 7 = 2$, chỉ có thể là $2 + 0$ (trùng 0) hoặc $1 + 1$ (hai chữ số giống nhau), loại.\\
- Nếu hàng nghìn là 6 thì hai chữ số còn lại có tổng bằng $9 - 6 = 3$, chọn hai chữ số là 2 và 1.\\
Do đó, số lớn nhất có thể lập được là $96\,210$.\\
\textbf{Đáp án của câu hỏi là B}
}
\end{ex}

\textbf{Phần 2. Trắc nghiệm đúng sai}
\setcounterpageref{ex}{0}

\begin{ex} %[4D1V01]
Xét tính đúng sai của các khẳng định sau về số tự nhiên:
\choiceTF
{\True Hiệu giữa số lớn nhất có 5 chữ số và số bé nhất có 5 chữ số là $89\,999$.}
{Số gồm $4$ chục nghìn, $15$ trăm và $6$ chục được viết là $40\,156$.}
{\True Khi làm tròn số $54\,500$ đến hàng nghìn ta được số $55\,000$.}
{\True Có đúng $9$ số có 5 chữ số giống nhau.}
\loigiai{
\begin{itemize}
    \item Số lớn nhất có 5 chữ số là $99\,999$, số bé nhất có 5 chữ số là $10\,000$. Hiệu là: $99\,999 - 10\,000 = 89\,999$. Khẳng định này Đúng.
    \item Ta có $4$ chục nghìn $= 40\,000$; $15$ trăm $= 1500$; $6$ chục $= 60$. Số đó là: $40\,000 + 1500 + 60 = 41\,560$ (không phải $40\,156$). Khẳng định này Sai.
    \item Chữ số hàng trăm là 5 ($5 \ge 5$) nên khi làm tròn đến hàng nghìn ta làm tròn lên thành $55\,000$. Khẳng định này Đúng.
    \item Các số có 5 chữ số giống nhau là: $11\,111; 22\,222; 33\,333; 44\,444; 55\,555; 66\,666; 77\,777; 88\,888; 99\,999$ (gồm đúng 9 số). Khẳng định này Đúng.
\end{itemize}
\begin{itemize}
    \item a) \textbf{Đúng}.
    \item b) \textbf{Sai}.
    \item c) \textbf{Đúng}.
    \item d) \textbf{Đúng}.
\end{itemize}
}
\end{ex}

\begin{ex} %[4D1V01]
Cho các số: $A = 58\,250$; $B = 58\,520$; $C = 58\,052$; $D = 58\,502$. Xét tính đúng sai của các khẳng định sau:
\choiceTF
{\True Thứ tự sắp xếp các số theo chiều giảm dần là: $B > D > A > C$.}
{\True Khi làm tròn cả 4 số trên đến hàng nghìn, có đúng 2 số được làm tròn thành $59\,000$.}
{\True Hiệu giữa số lớn nhất và số bé nhất trong bốn số trên bằng $468$.}
{Tổng của số lớn nhất và số bé nhất là một số chia hết cho $5$.}
\loigiai{
\begin{itemize}
    \item So sánh các số: $58\,520 > 58\,502 > 58\,250 > 58\,052$ tức là $B > D > A > C$. Khẳng định này Đúng.
    \item Làm tròn từng số đến hàng nghìn:
    \begin{itemize}
        \item $A = 58\,250$ có hàng trăm là 2 nên làm tròn thành $58\,000$.
        \item $B = 58\,520$ có hàng trăm là 5 nên làm tròn thành $59\,000$.
        \item $C = 58\,052$ có hàng trăm là 0 nên làm tròn thành $58\,000$.
        \item $D = 58\,502$ có hàng trăm là 5 nên làm tròn thành $59\,000$.
    \end{itemize}
    Như vậy có đúng 2 số ($B$ và $D$) làm tròn thành $59\,000$. Khẳng định này Đúng.
    \item Số lớn nhất là $B = 58\,520$, số bé nhất là $C = 58\,052$. Hiệu giữa hai số là: $58\,520 - 58\,052 = 468$. Khẳng định này Đúng.
    \item Tổng của số lớn nhất và số bé nhất là: $58\,520 + 58\,052 = 116\,572$. Chữ số tận cùng là 2 nên tổng này không chia hết cho 5. Khẳng định này Sai.
\end{itemize}
\begin{itemize}
    \item a) \textbf{Đúng}.
    \item b) \textbf{Đúng}.
    \item c) \textbf{Đúng}.
    \item d) \textbf{Sai}.
\end{itemize}
}
\end{ex}

\textbf{Phần 3. Trả lời ngắn}
\setcounterpageref{ex}{0}

\begin{ex} %[4D1V01]
Tìm số liền sau của số chẵn bé nhất có 5 chữ số khác nhau.
\shortans{10235}
\loigiai{
\textbf{Đáp số:} $10235$\\
Số chẵn bé nhất có 5 chữ số khác nhau là $10\,234$ (chữ số hàng chục nghìn là 1, hàng nghìn là 0, hàng trăm là 2, hàng chục là 3 và hàng đơn vị chẵn là 4).\\
Số liền sau của số $10\,234$ là:
$$10\,234 + 1 = 10\,235$$
}
\end{ex}

\begin{ex} %[4D1V01]
Một số tự nhiên có 5 chữ số có chữ số hàng đơn vị là $7$. Biết rằng khi xóa chữ số $7$ ở hàng đơn vị thì ta được số mới bé hơn số ban đầu $31\,849$ đơn vị. Tìm số ban đầu đó.
\shortans{35387}
\loigiai{
\textbf{Đáp số:} $35387$\\
Khi xóa chữ số 7 ở hàng đơn vị của một số thì số đó giảm đi 10 lần và 7 đơn vị, tức là số ban đầu bằng 10 lần số mới cộng thêm 7.\\
Hiệu giữa số ban đầu và số mới bằng 9 lần số mới cộng thêm 7 đơn vị:\\
$9 \times (\text{số mới}) + 7 = 31\,849$\\
$9 \times (\text{số mới}) = 31\,849 - 7 = 31\,842$\\
Số mới là: $31\,842 : 9 = 3538$.\\
Vậy số ban đầu là $35\,387$.
}
\end{ex}

\begin{ex} %[4D1V01]
Tìm số tự nhiên lớn nhất có 5 chữ số khác nhau mà khi làm tròn đến hàng nghìn ta được số $58\,000$.
\shortans{58497}
\loigiai{
\textbf{Đáp số:} $58497$\\
Để một số làm tròn đến hàng nghìn được $58\,000$ thì số đó phải nhỏ hơn $58\,500$ (tức chữ số hàng trăm phải từ 4 trở xuống: $0; 1; 2; 3; 4$).\\
Để số đó lớn nhất có 5 chữ số khác nhau:\\
- Chữ số hàng chục nghìn là 5.\\
- Chữ số hàng nghìn là 8.\\
- Chữ số hàng trăm lớn nhất có thể là 4.\\
- Chữ số hàng chục lớn nhất chưa dùng là 9.\\
- Chữ số hàng đơn vị lớn nhất chưa dùng (khác 5, 8, 4, 9) là 7.\\
Vậy số lớn nhất cần tìm là $58\,497$.
}
\end{ex}

\begin{ex} %[4D1V01]
Tính tổng của số lẻ lớn nhất có 5 chữ số khác nhau và số chẵn bé nhất có 5 chữ số khác nhau.
\shortans{108999}
\loigiai{
\textbf{Đáp số:} $108999$\\
- Số lẻ lớn nhất có 5 chữ số khác nhau là $98\,765$ (chữ số hàng chục nghìn là 9, hàng nghìn là 8, hàng trăm là 7, hàng chục là 6 và hàng đơn vị lẻ lớn nhất là 5).\\
- Số chẵn bé nhất có 5 chữ số khác nhau là $10\,234$.\\
Tổng của hai số đó là:
$$98\,765 + 10\,234 = 108\,999$$
}
\end{ex}

\begin{ex} %[4D1V01]
Có bao nhiêu số tự nhiên có 5 chữ số có chữ số tận cùng là $5$, lớn hơn $20\,000$ và bé hơn $25\,000$?
\shortans{500}
\loigiai{
\textbf{Đáp số:} $500$\\
Các số tự nhiên thỏa mãn đề bài tạo thành dãy số cách đều 10 đơn vị:
$$20\,005; 20\,015; 20\,025; \dots; 24\,995$$
Số lượng các số thỏa mãn là:
$$(24\,995 - 20\,005) : 10 + 1 = 4990 : 10 + 1 = 500\text{ (số)}$$
}
\end{ex}

\begin{ex} %[4D1V01]
Tìm số tự nhiên có 5 chữ số, biết rằng nếu viết thêm chữ số 2 vào bên phải số đó thì ta được số mới gấp 3 lần số nhận được khi viết thêm chữ số 2 vào bên trái số đó.
\shortans{85714}
\loigiai{
\textbf{Đáp số:} $85714$\\
Gọi số cần tìm là $x$ ($10\,000 \le x \le 99\,999$).\\
- Khi viết thêm chữ số 2 vào bên phải số đó, ta được số mới là: $10 \times x + 2$.\\
- Khi viết thêm chữ số 2 vào bên trái số đó, ta được số mới là: $200\,000 + x$.\\
Theo đề bài, ta có biểu thức:
\begin{align*}
10 \times x + 2 &= 3 \times (200\,000 + x) \\
10 \times x + 2 &= 600\,000 + 3 \times x \\
10 \times x - 3 \times x &= 600\,000 - 2 \\
7 \times x &= 599\,998 \\
x &= 599\,998 : 7 \\
x &= 85\,714
\end{align*}
Vậy số cần tìm là $85\,714$.
}
\end{ex}
